% =====================================================================
% 4. Aims, Objectives and Research Questions
% =====================================================================
\section{Aims, Objectives and Research Questions}
\label{sec:aims}

The aim of this project is to design, implement and evaluate a read-only dashboard that provides engineers, project leads and customers with an accurate and timely view of BMS tests, without interfering with the test rig.
The project is classified as hybrid, in that the system is developed as the means of answering four research questions.

\begin{description}[leftmargin=1.2em, labelindent=0pt, font=\normalfont\bfseries]
    \item[RQ1.] Can a read-only streaming pipeline with interchangeable data adapters deliver KPIs in near real time without measurably affecting the test rig?
    \item[RQ2.] Which filtering methods reduce measurement noise while preserving the visibility of genuine fault transients?
    \item[RQ3.] Can a tiered alert scheme with deduplication, hysteresis and shelving reduce the number of alerts without missing critical faults?
    \item[RQ4.] Do role-specific views reduce predicted task time and visual clutter while satisfying accessibility requirements?
\end{description}

The objectives are as follows:
\begin{enumerate}[label=\textbf{O\arabic*}]
    \item to derive a documented set of requirements from the partner's project brief;
    \item to implement simulator, replay and live-rig adapters behind a common interface;
    \item to generate test traces containing spikes, drifts and dropouts at known times;
    \item to implement and compare at least four filtering methods;
    \item to implement the alert engine and the three role-specific views;
    \item to conduct the four experiments described in Section~\ref{sec:evaluation} and to report the results against the criteria in Table~\ref{tab:success}.
\end{enumerate}

\subsection{Changes since the proposal}\label{subsec:changes-since-the-proposal}
RQ3 has been revised to measure the reduction in alerts directly, rather than by comparison with external industrial targets.
The original timeline also followed the partner's rollout schedule; as deployment, training and go-live are not assessed, the revised plan allocates dedicated time to the experiments (Table~\ref{tab:deliverables}).

\begin{table}[H]
    \centering\small
    \caption{Company and research deliverables. Only the research deliverables are assessed.}
    \label{tab:deliverables}
    \begin{tabularx}{\textwidth}{YY}
        \toprule
        \textbf{Company deliverable} & \textbf{Research deliverable} \\
        \midrule
        GUI deployed on the STT environment & Reference architecture and adapters \\
        Training, go-live and sign-off & Experiment scripts, datasets and results \\
        \bottomrule
    \end{tabularx}
\end{table}
